% !PRINTED 719
\weberpage{719}
%<<CONT

\gesperrt{Eine Einheit in $H_8$, die mit ihren Conjugirten von einerlei Zeichen ist,
ist, vom Zeichen abgesehen, das Quadrat einer Einheit.}

\begin{center}
\S.\,197.

\gesperrt{Recurrente Berechnung der Classenzahl im K\"orper $\Omega_m$,
wenn $m$ eine Potenz von 2 ist.}
\end{center}

Zur allgemeinen Bestimmung der Classenzahl im K\"orper $\Omega_m$ unter
der Voraussetzung, dass $m$ irgend eine Potenz von 2 und gr\"osser als 8
ist, wenden wir ein recurrirendes Verfahren an. Es sei also
\weq{1}{m = 2^\varkappa, \quad \mu = 2^{\varkappa-1}, \quad
\nu = 2^{\varkappa-2}, \quad \nu' = 2^{\varkappa-3}.}
Neben dem K\"orper $\Omega_m$ betrachten wir den K\"orper $\Omega_\mu$,
der ein Theiler von $\Omega_m$ ist, und den wir hier auch mit $\Omega_m'$
bezeichnen wollen. Es sei $h'$ die Classenzahl im K\"orper $\Omega_m'$,
d.\,h. die Zahl, die sich aus $h$ ergiebt, wenn $\varkappa$ in
$\varkappa - 1$ verwandelt wird, und wir setzen
\weq{2}{h = h' H:}
Setzen wir $h'$ als schon bekannt voraus, so kommt es nur noch auf die
Berechnung von $H$ an, was, wie sich zeigen wird, eine ganze Zahl ist.
Indem wir $h$ und $h'$ wie im \S.\,195 zerlegen, setzen wir
\[
\begin{aligned}
(3)\quad h &= k\,h_1, & \qquad h' &= k'\,h_1', \\
(4)\quad k &= k'A, & \qquad h_1 &= h_1'B,
\end{aligned}
\qquad H = AB,
\]
worin $k'$ aus $k$ und $h_1'$ aus $h_1$ durch die Vertauschung von
$\varkappa$ mit $\varkappa - 1$ hervorgeht. Wir wollen \"uberhaupt jetzt
durch Accente an den Buchstaben andeuten, dass $\varkappa - 1$ statt
$\varkappa$ gesetzt sein soll.

Nach \S.\,176, (7) ist dann
\[
\begin{aligned}
\Delta_1 &= 2^{(\varkappa-1)2^{\varkappa-2}-1}, \\
\Delta_1' &= 2^{(\varkappa-2)2^{\varkappa-3}-1},
\end{aligned}
\]
und folglich
\weq{5}{\Delta_1 = 2^{\varkappa\,2^{\varkappa-3}} \Delta_1'.}
Hat $E$ die Bedeutung \S.\,195, (12), und geht $E'$ aus $E$ hervor, wenn
$\varkappa$ durch $\varkappa - 1$ ersetzt wird, so setzen wir noch
\weq{6}{E = E'D,}
%CONT>>
